% ============================================================
%  VI. EXPERIMENTS
% ============================================================
\section{Experimental Setup}
\label{sec:experiments}

\subsection{Hardware Platform}

\ac{TERESA} is built on a Boston Dynamics Spot quadruped equipped with a
Unitree Z1 collaborative arm mounted on the rear payload deck. The primary
perception sensor is an Orbbec Femto Bolt \ac{RGBD} camera with an
\SI{87}{\degree} horizontal field of view, rigidly attached to the front of
the Spot chassis at a height of approximately \SI{0.5}{\metre} above the
ground. A secondary Intel RealSense D435 depth camera is mounted on the Z1
end-effector, providing close-range torso tracking during the approach and
scanning phases. All computation runs on an onboard Intel NUC 13 Pro
(i7-1360P) with an NVIDIA RTX 4060 Ti GPU for real-time inference. The
software stack is implemented in ROS 2 Humble, with the whole-body controller
running at \SI{100}{\hertz} via a \ac{QP} solver configured from the
parameters in \texttt{wbc\_params.yaml}.

\subsection{Search Time Experiment}

The first experiment quantifies the advantage of the Z1 arm during the
initial patient-search phase. In the \ac{TERESA} configuration, the arm
provides approximately \SI{300}{\degree} of angular coverage from each Spot
yaw position, allowing the \ac{EE}-mounted RealSense to observe a wide arc
without moving the base. In contrast, a Spot-only configuration with the
Orbbec camera fixed on the robot body must rely entirely on base rotation to
cover the full \SI{360}{\degree} azimuth.

We evaluate both configurations through a Monte Carlo simulation with
$N = \num{10000}$ trials. In each trial, the patient's yaw angle
$\theta$ is drawn uniformly from $[0\degree, 360\degree)$, representing an
unknown lying orientation. For the Spot+Arm configuration, the robot scans
from three base yaw positions ($0\degree$, $+30\degree$, $-30\degree$),
with the arm executing seven search poses per position (the
\texttt{search\_yaw\_angles} and \texttt{ik\_done} event count defined in
\texttt{wbc\_params.yaml}).  At each base yaw, the arm's
\SI{300}{\degree} angular coverage is modelled as a detection window
$[\psi-150\degree, \psi+150\degree]$ around the current heading $\psi$;
the patient is found when $\theta$ falls within this window, with detection
latency proportional to the angular offset from the window centre.  For the
Orbbec-Only configuration, the camera's effective \ac{FOV} of
\SI{77}{\degree} (the nominal \SI{87}{\degree} minus \SI{10}{\degree}
overlap for reliable detection) requires $N=\lceil 360/77\rceil = 5$
discrete base yaw positions for full \SI{360}{\degree} coverage. At each
position the base dwells for \SI{2.0}{\second} to acquire and process
frames, and rotates between positions at \SI{0.2}{\radian\per\second}.

The model parameters are traced directly to the real system configuration:
\SI{0.2}{\radian\per\second} base yaw velocity
(\texttt{search\_max\_angular\_vel}), \SI{1.2}{\second} arm pose timeout
(\texttt{search\_timeout\_per\_point}), and $\pm\SI{30}{\degree}$ search
yaw angles (\texttt{search\_yaw\_angles}).  Reported search times are
computed analytically from these parameters for each of the five canonical
patient yaw angles ($0\degree$, $45\degree$, $90\degree$, $135\degree$,
$180\degree$) and aggregated over the \num{10000} Monte Carlo trials.  The
full simulation source is available in
\texttt{scripts/search\_time\_model.py} with all parameters exposed as
configurable constants.

\subsection{Additional Experiments}

The following subsections describe the remaining experimental evaluations:
\ac{FAST} ultrasound probe positioning accuracy, body scanning completeness
and duration, and ablation studies of the perception-quality-aware velocity
controller.
